% =====================================================
% CAPA / ABERTURA
% =====================================================
\capaRoteiro{\faDesktop}{Slides-Guia de Projeção (8 Telas) \& Moodle}
  {Guia operacional da Frente 3 para as Oficinas de 120 minutos da SNCT 2026}
  {\chip{\faLayerGroup}{Extensão A2 (20\%)}\hspace{4pt}
   \chip{\faTools}{HTML5 · Bootstrap 5 · GitHub Pages · Moodle}\hspace{4pt}
   \chip{\faPlay}{Equipes de 3 estudantes}}

\vspace{0.4em}
\cartoesUC{Apresentação interativa de 8 telas em HTML/Bootstrap 5 publicada no GitHub Pages e integrada à sala do Moodle}

\begin{boxObjetivos}
\begin{itemize}
  \item \textbf{Construir os Slides-Guia de exatamente 8 telas em HTML5 + Bootstrap 5}, projetados para leitura a 6 metros de distância no telão do laboratório;
  \item \textbf{Sincronizar a numeração das telas com os passos do Micro-Roteiro} (Frente 1) e com os arquivos do Starter Kit (Frente 2);
  \item \textbf{Implementar navegação fluida por teclado} (setas $\leftarrow$ / $\rightarrow$) e contador visual de progresso com cronômetro de bloco;
  \item \textbf{Publicar os slides no GitHub Pages e incorporá-los na sala do Moodle institucional} (com papel de \textbf{Professor}), permitindo que cada aluno acompanhe também no seu monitor;
  \item \textbf{Exportar versão PDF offline} para o pen-drive de contingência da oficina.
\end{itemize}
\end{boxObjetivos}

\begin{mapaAula}
\etapaMapa{azulMedio}{1}{Missão desta semana}{Candidatura até 27/10 e divisão das 8 telas entre os 3 integrantes (Seção~\ref{sec:semana})}
\etapaMapa{verdeVivo}{2}{Roteiro das 8 telas}{O que entra em cada tela e como casar com os 120 minutos (Seção~\ref{sec:oito_telas})}
\etapaMapa{ambar}{3}{Código base HTML}{Template completo com navegação por teclado e fonte de telão (Seção~\ref{sec:codigo})}
\etapaMapa{roxo}{4}{Moodle e ensaio}{Publicação no GitHub Pages, link no Moodle e checklist (Seção~\ref{sec:marcos})}
\end{mapaAula}

% =====================================================
% 1. ONDE ESTA FRENTE SE ENCAIXA E TAREFAS DA SEMANA
% =====================================================
\section{Onde esta frente se encaixa e o que fazer nesta semana}\label{sec:semana}

Durante a oficina, quem está conduzindo no telão precisa ditar o ritmo da turma sem deixar ninguém para trás. Os \textbf{Slides-Guia de 8 Telas} funcionam como o painel de comando da oficina: mostram na tela grande qual arquivo abrir, qual linha editar e qual é o \emph{checkpoint} de vitória de cada bloco.

\begin{table}[H]
\centering
\small
\renewcommand{\arraystretch}{1.45}\rowcolors{2}{fundo}{white}
\begin{tabularx}{\textwidth}{RYY}
\cabecalho \cab{Frente de Trabalho} & \cab{Entregável Principal} & \cab{Conexão com a Frente de Slides} \\
1. Didática \& Micro-Roteiro & PDF de até 4 páginas + Sala no Moodle IFSC. & Cada passo do PDF cita a tela correspondente (\emph{``Veja a Tela 3 no telão''}). \\
2. Starter Kit \& GitHub & Repositório template + \texttt{starter-kit.zip}. & Fornece os trechos de código e o QR Code do projeto pronto para a Tela 1. \\
\textbf{3. Slides-Guia (Você)} & 8 telas em HTML no GitHub Pages + PDF no Moodle. & Guia visual da projeção e link de apoio na Seção 0 do Moodle. \\
4. Monitoria \& Logística & Escala de bancada, presença e satisfação. & Fornece o link/QR Code do questionário final para a Tela 8. \\
\end{tabularx}
\caption{Articulação entre a Frente 3 e as demais equipes da oficina.}
\end{table}

\subsection*{Pipeline Visual: Da Produção em Equipe à Aplicação nos Dias da SNCT}

O diagrama abaixo resume o fluxo completo de execução que conecta o trabalho desta semana ao dia da oficina presencial no laboratório:

\begin{center}
\begin{adjustbox}{max width=\linewidth}
\begin{tikzpicture}[
    node distance=0.45cm,
    etapa/.style={rectangle, rounded corners=6pt, draw=#1, fill=#1!8, line width=1pt,
                  text width=2.85cm, minimum height=2.35cm, align=center, inner sep=5pt, font=\scriptsize\sffamily},
    seta/.style={-{Latex[length=2.6mm, width=2mm]}, line width=1.2pt, color=azul}
]
\node[etapa=azulMedio] (f1) {
  {\bfseries\color{azulMedio}1. PRODUÇÃO}\\[2pt]
  \textbf{Frentes 1 a 4}\\
  Codificação do Starter Kit, Slides (8 telas), Roteiro A4 e Pesquisa.\\[2pt]
  {\tiny\color{suave}(Semana 28--30/10)}
};

\node[etapa=ambar, right=of f1] (f2) {
  {\bfseries\color{ambar!85!black}2. TESTE NO LAB}\\[2pt]
  \textbf{Revisão Cruzada}\\
  Teste real nas máquinas do lab, código 100\% igual ao PDF e Plano B.\\[2pt]
  {\tiny\color{suave}(Pré-publicação)}
};

\node[etapa=verdeVivo, right=of f2] (f3) {
  {\bfseries\color{verde}3. MOODLE IFSC}\\[2pt]
  \textbf{Alunos = Professores}\\
  Criação das 5 seções nas Salas do Curso 1 e do Curso 2 no Moodle.\\[2pt]
  {\tiny\color{suave}(Chave: SNCT2026)}
};

\node[etapa=ciano, right=of f3] (f4) {
  {\bfseries\color{ciano}4. KIT + CÓDIGOS}\\[2pt]
  \textbf{Disponibilização}\\
  Upload do \texttt{.zip}, PDF, link GitHub Template e Slides na sala.\\[2pt]
  {\tiny\color{suave}(Pronto p/ público)}
};

\node[etapa=roxo, right=of f4] (f5) {
  {\bfseries\color{roxo}5. DIAS DA SNCT}\\[2pt]
  \textbf{Aplicação (120 min)}\\
  Oficina nas bancadas, site no celular (QR Code) e avaliação final.\\[2pt]
  {\tiny\color{suave}(Evento + Relatório)}
};

\draw[seta] (f1.east) -- (f2.west);
\draw[seta] (f2.east) -- (f3.west);
\draw[seta] (f3.east) -- (f4.west);
\draw[seta] (f4.east) -- (f5.west);
\end{tikzpicture}
\end{adjustbox}
\end{center}


\begin{atencao}{Ação imediata para esta semana (ausência docente em sala)}
\begin{enumerate}[leftmargin=1.4em, itemsep=0.25em]
  \item \textbf{Até amanhã (terça-feira, 27/10, às 23h59):} acesse o painel da disciplina com seu login Google e candidate-se na frente \emph{Slides-Guia de Projeção (8 Telas)} do Curso 1 ou do Curso 2.
  \item \textbf{Durante a aula desta semana:} dividam as 8 telas entre os 3 integrantes conforme a tabela abaixo, publiquem a versão inicial no GitHub Pages e registrem o link no painel.
\end{enumerate}
\end{atencao}

\begin{table}[H]
\centering
\small
\renewcommand{\arraystretch}{1.45}\rowcolors{2}{fundo}{white}
\begin{tabularx}{\textwidth}{RYY}
\cabecalho \cab{Papel na Equipe} & \cab{Responsabilidade Técnica} & \cab{Entrega até o fim desta semana} \\
Membro A — Estrutura \& Telas 1 a 3 & Montar o template base HTML/CSS/JS e redigir as Telas 1 (Boas-vindas), 2 (Moodle) e 3 (Bloco 1). & Navegação por setas funcionando e Telas 1--3 prontas. \\
Membro B — Telas 4 a 6 (Mão na Massa) & Redigir as Telas 4 e 5 (Bloco 2: código central) e a Tela 6 (Bloco 3: GitHub Pages). & Blocos de código com destaque visual e alto contraste. \\
Membro C — Telas 7 e 8, Deploy \& Moodle & Redigir as Telas 7 (QR Code mobile) e 8 (Avaliação), ativar GitHub Pages e linkar no Moodle. & URL pública ativa e inserida na Seção 0 do Moodle. \\
\end{tabularx}
\caption{Divisão de trabalho recomendada para os 3 estudantes da Frente 3.}
\end{table}

% =====================================================
% 2. MAPA DAS 8 TELAS DOS SLIDES-GUIA
% =====================================================
\section{Mapa Cirúrgico das 8 Telas (120 Minutos)}\label{sec:oito_telas}

Regra de ouro para projeção em laboratório: \textbf{nunca coloque mais de 12 linhas de código em uma tela} e use tamanho mínimo de \texttt{1.35rem} (22px) para o texto.

\begin{table}[H]
\centering
\small
\renewcommand{\arraystretch}{1.45}\rowcolors{2}{fundo}{white}
\begin{tabularx}{\textwidth}{RYY}
\cabecalho \cab{Tela / Tempo} & \cab{Equipe Curso 1: Meu Primeiro Site no Ar} & \cab{Equipe Curso 2: Painel Verde (Coleta \& Ecopontos)} \\
Tela 1 (00--10 min) & QR Code grande da página pessoal modelo publicada. & QR Code grande do Painel Verde de Garopaba publicado. \\
Tela 2 (10--15 min) & Acesso à \textbf{Sala Curso 1} no Moodle + chave \texttt{SNCT2026}. & Acesso à \textbf{Sala Curso 2} no Moodle + chave \texttt{SNCT2026}. \\
Tela 3 (15--35 min) & Bloco 1: Baixar \texttt{starter-kit-snct.zip} e editar \texttt{:root}. & Bloco 1: Baixar \texttt{starter-kit-painel-verde.zip} e editar título/cores. \\
Tela 4 (35--55 min) & Bloco 2A: Editar os 3 cartões de destaques em \texttt{index.html}. & Bloco 2A: Cadastrar novo ecoponto no array \texttt{LOCAIS\_COLETA} (\texttt{js/app.js}). \\
Tela 5 (55--70 min) & Bloco 2B: Configurar o link direto de WhatsApp (\texttt{wa.me}). & Bloco 2B: Testar os botões de filtro (\texttt{filtrarEcopontos}) e link do Google Maps. \\
Tela 6 (70--95 min) & Bloco 3: Upload no GitHub e ativação do GitHub Pages. & Bloco 3: Upload do Painel Verde no GitHub e ativação do GitHub Pages. \\
Tela 7 (95--110 min) & Bloco 4: Gerar QR Code da URL \texttt{.github.io} e postar no Mural. & Bloco 4: Gerar QR Code do Painel Verde e postar no Mural do Curso 2. \\
Tela 8 (110--120 min) & QR Code da Pesquisa de Satisfação na Seção 4 do Moodle (Curso 1). & QR Code da Pesquisa de Satisfação na Seção 4 do Moodle (Curso 2). \\
\end{tabularx}
\caption{Roteiro das 8 telas separado para a Equipe do Curso 1 e para a Equipe do Curso 2.}
\end{table}

% =====================================================
% 3. CÓDIGO BASE COMPLETO DOS SLIDES-GUIA
% =====================================================
\section{Código Base Completo em HTML5 + Bootstrap 5}\label{sec:codigo}

\passo{3.1}{Criar o arquivo \texttt{index.html} dos Slides-Guia}

O código abaixo é 100\% funcional (regra \emph{Zero-Ellipsis}): suporta setas do teclado ($\leftarrow$ / $\rightarrow$), botões na tela e barra de progresso. Basta duplicar a estrutura \texttt{<section class="slide">...} para completar as 8 telas:

\begin{boxCodigoHTML}[slides-snct/index.html]
<!DOCTYPE html>
<html lang="pt-BR">
<head>
  <meta charset="UTF-8">
  <meta name="viewport" content="width=device-width, initial-scale=1.0">
  <title>Slides-Guia — Oficina SNCT 2026 IFSC Garopaba</title>
  <link href="https://cdn.jsdelivr.net/npm/bootstrap@5.3.3/dist/css/bootstrap.min.css"
        rel="stylesheet">
  <style>
    body { background: #0f172a; color: #f8fafc; font-size: 1.4rem; overflow: hidden; }
    .slide { display: none; height: calc(100vh - 76px); padding: 48px 64px; }
    .slide.ativo { display: flex; flex-direction: column; justify-content: center; }
    .selo-bloco { background: #16a34a; color: #fff; padding: 6px 16px;
                  border-radius: 999px; font-weight: 700; font-size: 1rem; }
    pre.codigo-telao { background: #1e293b; color: #38bdf8; padding: 24px;
                       border-radius: 12px; font-size: 1.25rem; max-height: 420px; }
    .barra-rodape { height: 76px; background: #1e293b; padding: 0 32px;
                    display: flex; align-items: center; justify-content: space-between; }
  </style>
</head>
<body>
  <!-- TELA 1 -->
  <section class="slide ativo">
    <div><span class="selo-bloco">TELA 1 DE 8 · ACOLHIDA (00–10 MIN)</span></div>
    <h1 class="display-4 fw-bold mt-3">Meu Primeiro Site no Ar em 120 Minutos</h1>
    <p class="text-info fs-3">IFSC Campus Garopaba — Semana Nacional de Ciencia e Tecnologia</p>
    <div class="p-4 mt-3 rounded-4 bg-dark border border-success">
      <strong>Meta de hoje:</strong> voce vai sair deste laboratorio com seu proprio site
      publicado na internet e abrindo no celular via QR Code!
    </div>
  </section>

  <!-- TELA 2 -->
  <section class="slide">
    <div><span class="selo-bloco">TELA 2 DE 8 · SALA NO MOODLE (10–15 MIN)</span></div>
    <h2 class="display-5 fw-bold mt-3">Passo Inicial: Abra a Sala no Moodle IFSC</h2>
    <pre class="codigo-telao mt-3">1. Acesse a sala da Oficina SNCT no Moodle do IFSC
2. Chave de inscricao da turma: SNCT2026
3. Na Secao 0, abra o Micro-Roteiro PDF
4. Na Secao 1, clique em starter-kit-snct.zip</pre>
  </section>

  <!-- TELA 3 -->
  <section class="slide">
    <div><span class="selo-bloco">TELA 3 DE 8 · BLOCO 1 (15–35 MIN)</span></div>
    <h2 class="display-5 fw-bold mt-3">Personalizando Cores no arquivo css/estilo.css</h2>
    <pre class="codigo-telao mt-3">:root {
  --cor-primaria: #2563eb;  /* Troque pela sua cor favorita */
  --cor-fundo:    #eff6ff;
}</pre>
  </section>

  <footer class="barra-rodape">
    <span id="indicadorTela" class="fw-bold">Tela 1 / 3</span>
    <div class="d-flex gap-2">
      <button class="btn btn-outline-light btn-lg" onclick="mudarSlide(-1)">Anterior</button>
      <button class="btn btn-success btn-lg" onclick="mudarSlide(1)">Proxima Tela</button>
    </div>
  </footer>

  <script>
    let indiceAtual = 0;
    const slides = document.querySelectorAll('.slide');
    const indicador = document.getElementById('indicadorTela');
    function mudarSlide(delta) {
      slides[indiceAtual].classList.remove('ativo');
      indiceAtual = Math.max(0, Math.min(slides.length - 1, indiceAtual + delta));
      slides[indiceAtual].classList.add('ativo');
      indicador.textContent = `Tela ${indiceAtual + 1} / ${slides.length}`;
    }
    document.addEventListener('keydown', (ev) => {
      if (ev.key === 'ArrowRight') mudarSlide(1);
      if (ev.key === 'ArrowLeft') mudarSlide(-1);
    });
  </script>
</body>
</html>
\end{boxCodigoHTML}

% =====================================================
% 4. PUBLICAÇÃO NO MOODLE, MARCOS E CHECKLIST
% =====================================================
\section{Publicação no Moodle, Marcos e Checklist de Entrega}\label{sec:marcos}

\passo{4.1}{Publicar no GitHub Pages e adicionar na Seção 0 do Moodle}
\begin{enumerate}[leftmargin=1.4em, itemsep=0.25em]
  \item Ative o GitHub Pages no repositório dos slides (\botao{Settings}\seta\botao{Pages}\seta Branch \texttt{main});
  \item Acesse a sala do curso no Moodle institucional (com seu perfil de \textbf{Professor}), clique em \botao{Adicionar atividade ou recurso}\seta\botao{URL} na \textbf{Seção 0} e cole o link público dos slides;
  \item Exporte também uma cópia em PDF (\tecla{Ctrl}+\tecla{P}\seta\emph{Salvar como PDF}) e anexe na Seção 0 como contingência offline.
\end{enumerate}

\begin{checklist}[Checklist de prontidão da Frente 3 (Slides-Guia \& Moodle)]
\begin{itemize}[label={\caixaMarcar}, itemsep=0.35em]
  \item Os 3 estudantes registraram candidatura na frente até 27/10 às 23h59 no painel da disciplina;
  \item A apresentação possui exatamente as 8 telas mapeadas na Seção~\ref{sec:oito_telas}, com navegação por teclado;
  \item Os números das telas e os nomes de arquivos coincidem 100\% com o Micro-Roteiro da Frente 1;
  \item Todos os blocos de código têm fonte legível do fundo do laboratório ($\ge$ 20px) e alto contraste;
  \item Os slides estão publicados no GitHub Pages e linkados na Seção 0 da sala do Moodle institucional;
  \item Uma cópia em PDF dos slides foi salva nos pen-drives do Plano B e o link foi registrado no painel.
\end{itemize}
\end{checklist}

\vfill
{\centering\footnotesize\color{suave}\faMapMarkerAlt\hspace{0.4em}IFSC Câmpus Garopaba — Programação para a Internet 2 (Curricularização da Extensão SNCT 2026)\par}

